What we observed are below,
1. We refer the "history" information as the results (or intermediate results) from previous computation. For example, if we run a SSSP from s to t in the previous time snapshot. Then, the graph udpates, and if we need to run a SSSP from s to t again, we can reuse many of the results from previous computation.
2. Most dynamic graph computation systems do not actually use this type of history information. Given a graph update, they would quickly integrate the update to the graph representation. Then, the graph query (e.g., SSSP) will simply be rerun on the updated graph.
3. There are also some works overly rely on the history computation (or query-reliant). For example, https://ieeexplore.ieee.org/abstract/document/9442368, this paper from [TPDS '21] would do the following. After a SSSP run, it would record the SSSP tree. Then, after the graph updated, it would locate the changed subgraphs and only recompute that part. This method is smart, but it is impractical because it depends on the source (or root) node s. If s changes, the SSSP tree would be different.
4. So, we think, the current solutions are on the two extreme, one is not using the previous computation, the other is query-dependent. Motivated by this, we design this new Graph IR, which updates a few essential fundamental graph algorithms when the graph updates. Then, for new graph queries, we would compute them based on the Graph IR. This would make our design query-agnostic.